\documentclass[12pt]{article}
\usepackage[utf8]{inputenc}
\usepackage[T1]{fontenc}
\usepackage[brazil]{babel}
\usepackage[a4paper,margin=2.5cm]{geometry}
\usepackage{enumitem}
\setlist[itemize]{noitemsep,topsep=4pt}
\title{Sprint 3 - S3\_08\\GitHub, Pull Requests e Integração Contínua}
\author{Matheus Moreli Durão Silva\\Matrícula 12621CBS217\\Cibersegurança - Bacharelado - Noturno}
\date{2026}
\begin{document}
\maketitle

\section*{Entregável}
Histórico Git, Pull Request e GitHub Actions. Este documento registra o objetivo do card, a solução desenvolvida e a forma de conferência. Os arquivos de código e as evidências citadas permanecem na mesma pasta deste documento.

\section*{O que foi estudado e aplicado}
\begin{itemize}
\item Branches e commits organizados
\item Pull Request documentado
\item Workflow que executa testes
\item Workflow que executa verificação de estilo
\end{itemize}

\section*{Evidência e limites}
A Pull Request \#3 do projeto principal foi criada a partir da branch \texttt{melhoria-trabalho-2-completo}. O GitHub Actions executou \texttt{testes-e-estilo} com sucesso na execução 38472278580. Links e detalhes estão em \texttt{evidencia\_actions.md}. Revisão por colega e resolução de conflito nesta PR ainda não foram realizadas. O workflow testa Python e estilo; a prova Docker de 24 horas permanece registrada na issue \#4.

\section*{Conclusão}
O artefato foi organizado para demonstrar tanto o conceito estudado quanto sua aplicação prática. As evidências de execução só são consideradas válidas quando produzidas pela ferramenta correspondente.
\end{document}
